\chapter{Что машина делает ночью}
% полная версия: chapters/20_sleep.tex

Сон --- не выключение машины, а смена режима. Это видно по радиостанциям. В глубоком сне приглушены почти все: ацетилхолин, норадреналин, серотонин. Во сне со сновидениями ацетилхолин возвращается к дневному уровню, норадреналин и серотонин почти молчат, а выход на мышцы отключён тормозами --- поэтому во сне мы не бегаем за тем, от чего убегаем во сне. Всё это измерено.

\section*{Когда пора спать}

Представьте конденсатор, который заряжается весь день, пока вы бодрствуете, и разряжается ночью. Заряд --- это «давление сна». Засыпаете вы, когда заряд достигает верхнего порога, просыпаетесь, когда он опускается до нижнего. А сами пороги плавно ходят вверх и вниз вместе с суточным ритмом из главы 16. Такая простая схема сама приходит к шестнадцати часам бодрствования и восьми часам сна. И объясняет, почему после бессонной ночи мы спим не вдвое дольше, а глубже: полный конденсатор разряжается быстрее. Кандидат на роль «заряда» --- вещество, которое накапливается в мозге при работе; что это именно оно, пока гипотеза.

\pencilfig{120}{%
% figures/sleep_{S,H,L}.dat from models/sleep_models.py, t in hours
\begin{scope}[x=0.1cm, y=2.6cm]
\foreach \a/\b in {0/4.3,21.2/29.3,45.4/53.4,69.5/72}{\fill[pcBlue, opacity=0.1] (\a,0) rectangle (\b,1);}
\draw[pen] (0,0) -- (72,0);
\draw[penlite=pcRed, densely dashed] plot file {figures/sleep_H.dat};
\draw[penlite=pcGreen, densely dashed] plot file {figures/sleep_L.dat};
\draw[pcBlue, line width=1.0pt] plot file {figures/sleep_S.dat};
\foreach \t/\l in {0/0,24/сутки,48/двое,72/трое}{\draw[pcGray, line width=0.5pt] (\t,0) -- (\t,-0.04); \node[lbl, anchor=base] at (\t,-0.16) {\l};}
\node[lbl, text=pcBlue] at (25.25,0.93) {сон}; \node[lbl, text=pcBlue] at (49.4,0.93) {сон};
\end{scope}
\node[lbl, text=pcRed, anchor=west] at (7.3,2.0) {заснуть};
\node[lbl, text=pcGreen!70!black, anchor=west] at (7.3,0.62) {проснуться};
\node[lbl, text=pcBlue, anchor=west] at (7.3,1.35) {давление сна};
}{Давление сна копится днём и тает ночью; засыпаем у верхнего порога, просыпаемся у нижнего, а оба порога качаются вместе с сутками.}

\section*{Медленные качели}

В глубоком сне целые области мозга синхронно то включаются, то выключаются --- реже раза в секунду. Это уже знакомая схема: группа, которая возбуждает сама себя и устаёт, --- генератор. Днём ацетилхолин подавляет усталость нейронов, и та же группа работает ровно; ночью его мало, и сеть начинает качаться.

\pencilfig{20}{\pencilart{20}}{Ночью сеть качается медленно: то почти все нейроны работают, то почти все молчат.}

\section*{Перемотка вперёд}

В этих качелях происходит самое интересное. Цепочки срабатываний нейронов, случившиеся днём, «проигрываются» заново --- и в разы быстрее: в отделе памяти примерно в двадцать раз, в коре в несколько раз. Если искусственно усилить согласованность этих проигрываний с медленными качелями, память улучшается --- это причинный опыт. Зачем это нужно? Правдоподобная гипотеза: медленная долговременная память учится на повторах вперемешку со старым, чтобы новое не затёрло прежнее. Инженеры знают эту беду: нейросеть, обученная на новой задаче, забывает старую, если не повторять старые примеры.

\section*{Подрезка связей}

Ещё одна измеренная вещь: после сна контакты между нейронами в среднем становятся меньше --- на единицы-десятки процентов, в пропорции к своему размеру. Как если бы все громкости немного убавили одним движением: соотношения сохранились, выученное осталось, а место для завтрашнего обучения освободилось. Что это главная задача сна --- гипотеза.

\section*{Сновидения и уборка}

Сон со сновидениями похож на стендовое испытание: машина работает на полную мощность, но входы заменены внутренними сигналами, а выход отключён. Зачем это нужно, мы не знаем. Идея, что во сне мозг «промывается» от отходов, пока спорна: есть опыты и за, и против.

Самое удивительное: сон бывает местным. После долгого бодрствования отдельные участки коры бодрствующего человека на мгновение «засыпают», замолкая посреди работы.

\fullversion{20}
